% GENERATED FILE. Edit canonical lessons, then run npm run generate:books.
% canonical-source-hash: fnv1a64:2e4fb758d6063732
% canonical-lessons: KA-C352-dange, KA-C352-kranti, KA-C352-gulama, KA-C352-bidugade-madu, KA-C352-vasapadisiko

\chapter{A Revolt Against the King}
\label{ch:ka-a-revolt-against-the-king}
\hlchaptermodality{352}

\begin{chapteropening}
\textbf{By the end of this chapter:} \emph{I can say a revolt, a revolution, a slave, to set free and to capture.}

Complete the last lesson of chapter 352: I can say a revolt, a revolution, a slave, to set free and to capture.
\end{chapteropening}

\section[\texorpdfstring{\kn{ದಂಗೆ} (daṅge)}{ದಂಗೆ (daṅge)}]{\kn{ದಂಗೆ} (daṅge) — a revolt; an uprising}

\label{lesson:KA-C352-dange}



\begin{quote}
Before the new one: say the Kannada for to engulf, then the Kannada for to be washed away.
\end{quote}

\subsection*{You'll want to know: \kn{ದಂಗೆ}}
\textbf{\kn{ದಂಗೆ}} — \emph{daṅge} — \textquotedblleft{}a revolt; an uprising\textquotedblright{}.

\begin{quote}
\kn{ಜನರು} \kn{ಆ} \kn{ರಾಜನ} \kn{ವಿರುದ್ಧ} \kn{ದಂಗೆ} \kn{ಎದ್ದರು}.
\end{quote}

\begin{quote}
\emph{janaru ā rājana viruddha daṅge eddaru.}
\end{quote}

\begin{quote}
The people rose in revolt against that king.
\end{quote}

One \textquotedblleft{}rises\textquotedblright{} in revolt: \kn{ದಂಗೆ} \kn{ಏಳು}.

\begin{cousinweb}
\textbf{\kn{ದಂಗೆ}} (\emph{daṅge}) is usually traced to Hindustani \emph{daṅgā}, \textquotedblleft{}a riot, a disturbance\textquotedblright{}; its earlier history is not settled.
\end{cousinweb}

\begin{grammarlens}[title={What you've built}]
One more word for telling what happened.
\end{grammarlens}

\subsection*{Your turn}
\begin{itemize}
\raggedright
  \item \emph{Say it:} \emph{daṅge}
  \item \emph{Say it:} \emph{daṅge}, once more
  \item \emph{Say it:} \emph{janaru ā rājana viruddha daṅge eddaru.} — The people rose in revolt against that king.
  \item \emph{Recall:} say the Kannada for to engulf, then the Kannada for to be washed away, then say \emph{daṅge} again
\end{itemize}

\begin{tcolorbox}[breakable,colback=teal!4,colframe=teal!35!black,title={Before you move on}]
What does \kn{ದಂಗೆ} mean? (\textquotedblleft{}A revolt; an uprising\textquotedblright{}.) Say it once more.
\end{tcolorbox}
\section[\texorpdfstring{\kn{ಕ್ರಾಂತಿ} (krānti)}{ಕ್ರಾಂತಿ (krānti)}]{\kn{ಕ್ರಾಂತಿ} (krānti) — a revolution}

\label{lesson:KA-C352-kranti}



\begin{quote}
Before the new one: say the Kannada for to be washed away, then the Kannada for a revolt.
\end{quote}

\subsection*{You'll want to know: \kn{ಕ್ರಾಂತಿ}}
\textbf{\kn{ಕ್ರಾಂತಿ}} — \emph{krānti} — \textquotedblleft{}a revolution\textquotedblright{}.

\begin{quote}
\kn{ಆ} \kn{ಕ್ರಾಂತಿ} \kn{ದೇಶದ} \kn{ಇತಿಹಾಸವನ್ನೇ} \kn{ಬದಲಿಸಿತು}.
\end{quote}

\begin{quote}
\emph{ā krānti dēśada itihāsavannē badalisitu.}
\end{quote}

\begin{quote}
That revolution changed the country's history.
\end{quote}

Plural: \kn{ಕ್ರಾಂತಿಗಳು}, \emph{krāntigaḷu}.

Also a great change: \kn{ಹಸಿರು} \kn{ಕ್ರಾಂತಿ}, the green revolution in farming.

\begin{cousinweb}
\textbf{\kn{ಕ್ರಾಂತಿ}} (\emph{krānti}) is borrowed from Sanskrit \emph{krānti}, \textquotedblleft{}going, proceeding, a step\textquotedblright{}, and \textquotedblleft{}overcoming\textquotedblright{} (Monier-Williams); \textquotedblleft{}a revolution\textquotedblright{} is a modern sense.
\end{cousinweb}

\begin{grammarlens}[title={What you've built}]
One more word for telling what happened.
\end{grammarlens}

\subsection*{Your turn}
\begin{itemize}
\raggedright
  \item \emph{Say it:} \emph{krānti}
  \item \emph{Say it:} \emph{krānti}, once more
  \item \emph{Say it:} \emph{ā krānti dēśada itihāsavannē badalisitu.} — That revolution changed the country's history.
  \item \emph{Recall:} say the Kannada for to be washed away, then the Kannada for a revolt, then say \emph{krānti} again
\end{itemize}

\begin{tcolorbox}[breakable,colback=teal!4,colframe=teal!35!black,title={Before you move on}]
What does \kn{ಕ್ರಾಂತಿ} mean? (\textquotedblleft{}A revolution\textquotedblright{}.) Say it once more.
\end{tcolorbox}
\section[\texorpdfstring{\kn{ಗುಲಾಮ} (gulāma)}{ಗುಲಾಮ (gulāma)}]{\kn{ಗುಲಾಮ} (gulāma) — a slave}

\label{lesson:KA-C352-gulama}



\begin{quote}
Before the new one: say the Kannada for a revolt, then the Kannada for a revolution.
\end{quote}

\subsection*{You'll want to know: \kn{ಗುಲಾಮ}}
\textbf{\kn{ಗುಲಾಮ}} — \emph{gulāma} — \textquotedblleft{}a slave\textquotedblright{}.

\begin{quote}
\kn{ಕಥೆಯಲ್ಲಿ} \kn{ಆ} \kn{ಗುಲಾಮ} \kn{ಕೊನೆಗೆ} \kn{ಸ್ವತಂತ್ರನಾದನು}.
\end{quote}

\begin{quote}
\emph{katheyalli ā gulāma konege svatantranādanu.}
\end{quote}

\begin{quote}
In the story the slave became free at last.
\end{quote}

Plural: \kn{ಗುಲಾಮರು}, \emph{gulāmaru}.

\kn{ಗುಲಾಮಗಿರಿ} is slavery.

\begin{cousinweb}
\textbf{\kn{ಗುಲಾಮ}} (\emph{gulāma}) is a loan from Arabic \emph{ġulām}, \textquotedblleft{}a boy, a servant, a slave\textquotedblright{}, which came through Persian and Hindustani.
\end{cousinweb}

\begin{grammarlens}[title={What you've built}]
One more word for telling what happened.
\end{grammarlens}

\subsection*{Your turn}
\begin{itemize}
\raggedright
  \item \emph{Say it:} \emph{gulāma}
  \item \emph{Say it:} \emph{gulāma}, once more
  \item \emph{Say it:} \emph{katheyalli ā gulāma konege svatantranādanu.} — In the story the slave became free at last.
  \item \emph{Recall:} say the Kannada for a revolt, then the Kannada for a revolution, then say \emph{gulāma} again
\end{itemize}

\begin{tcolorbox}[breakable,colback=teal!4,colframe=teal!35!black,title={Before you move on}]
What does \kn{ಗುಲಾಮ} mean? (\textquotedblleft{}A slave\textquotedblright{}.) Say it once more.
\end{tcolorbox}
\section[\texorpdfstring{\kn{ಬಿಡುಗಡೆ} \kn{ಮಾಡು} (biḍugaḍe māḍu)}{ಬಿಡುಗಡೆ ಮಾಡು (biḍugaḍe māḍu)}]{\kn{ಬಿಡುಗಡೆ} \kn{ಮಾಡು} (biḍugaḍe māḍu) — to set free; to release}

\label{lesson:KA-C352-bidugade-madu}



\begin{quote}
Before the new one: say the Kannada for a revolution, then the Kannada for a slave.
\end{quote}

\subsection*{You'll want to know: \kn{ಬಿಡುಗಡೆ} \kn{ಮಾಡು}}
\textbf{\kn{ಬಿಡುಗಡೆ} \kn{ಮಾಡು}} — \emph{biḍugaḍe māḍu} — \textquotedblleft{}to set free; to release\textquotedblright{}.

\begin{quote}
\kn{ಹೊಸ} \kn{ರಾಜ} \kn{ಎಲ್ಲ} \kn{ಗುಲಾಮರನ್ನು} \kn{ಬಿಡುಗಡೆ} \kn{ಮಾಡಿದನು}.
\end{quote}

\begin{quote}
\emph{hosa rāja ella gulāmarannu biḍugaḍe māḍidanu.}
\end{quote}

\begin{quote}
The new king set all the slaves free.
\end{quote}

Past (I …): \kn{ಬಿಡುಗಡೆ} \kn{ಮಾಡಿದೆ}, \emph{biḍugaḍe māḍide}.

Also used of a film or a book coming out: \kn{ಸಿನಿಮಾ} \kn{ಬಿಡುಗಡೆ} \kn{ಆಯಿತು}.

\begin{cousinweb}
\textbf{\kn{ಬಿಡುಗಡೆ} \kn{ಮಾಡು}} (\emph{biḍugaḍe māḍu}) is native Dravidian: \kn{ಬಿಡುಗಡೆ}, \textquotedblleft{}release\textquotedblright{}, is built on \emph{biḍu}, \textquotedblleft{}to let loose, to let go\textquotedblright{} (DEDR 5393), with \emph{māḍu}, \textquotedblleft{}to do\textquotedblright{} (DEDR 4797).
\end{cousinweb}

\begin{grammarlens}[title={What you've built}]
One more word for telling what happened.
\end{grammarlens}

\subsection*{Your turn}
\begin{itemize}
\raggedright
  \item \emph{Say it:} \emph{biḍugaḍe māḍu}
  \item \emph{Say it:} \emph{biḍugaḍe māḍu}, once more
  \item \emph{Say it:} \emph{hosa rāja ella gulāmarannu biḍugaḍe māḍidanu.} — The new king set all the slaves free.
  \item \emph{Recall:} say the Kannada for a revolution, then the Kannada for a slave, then say \emph{biḍugaḍe māḍu} again
\end{itemize}

\begin{tcolorbox}[breakable,colback=teal!4,colframe=teal!35!black,title={Before you move on}]
What does \kn{ಬಿಡುಗಡೆ} \kn{ಮಾಡು} mean? (\textquotedblleft{}To set free; to release\textquotedblright{}.) Say it once more.
\end{tcolorbox}
\section[\texorpdfstring{\kn{ವಶಪಡಿಸಿಕೊ} (vaśapaḍisiko)}{ವಶಪಡಿಸಿಕೊ (vaśapaḍisiko)}]{\kn{ವಶಪಡಿಸಿಕೊ} (vaśapaḍisiko) — to capture; to take control of}

\label{lesson:KA-C352-vasapadisiko}



\begin{quote}
Before the new one: say the Kannada for a slave, then the Kannada for to set free.
\end{quote}

\subsection*{You'll want to know: \kn{ವಶಪಡಿಸಿಕೊ}}
\textbf{\kn{ವಶಪಡಿಸಿಕೊ}} — \emph{vaśapaḍisiko} — \textquotedblleft{}to capture; to take control of\textquotedblright{}.

\begin{quote}
\kn{ದಂಗೆಕೋರರು} \kn{ಕೋಟೆಯನ್ನು} \kn{ವಶಪಡಿಸಿಕೊಂಡರು}.
\end{quote}

\begin{quote}
\emph{daṅgekōraru kōṭeyannu vaśapaḍisikoṇḍaru.}
\end{quote}

\begin{quote}
The rebels captured the fort.
\end{quote}

Past (I …): \kn{ವಶಪಡಿಸಿಕೊಂಡೆ}, \emph{vaśapaḍisikoṇḍe}.

\kn{ವಶ}, control, with -\kn{ಪಡಿಸು}, to bring, and -\kn{ಕೊ}, for oneself. The police \kn{ವಶಪಡಿಸಿಕೊ} stolen goods too.

\begin{cousinweb}
\textbf{\kn{ವಶಪಡಿಸಿಕೊ}} (\emph{vaśapaḍisiko}) is built from Sanskrit \emph{vaśa}, \textquotedblleft{}will, authority, power\textquotedblright{} (CDIAL 11430), with the native verb-makers -\emph{paḍisu} (DEDR 3853) and -\emph{ko}, from \emph{koḷ}, \textquotedblleft{}to take\textquotedblright{} (DEDR 2151).
\end{cousinweb}

\begin{grammarlens}[title={What you've built}]
One more word for telling what happened.
\end{grammarlens}

\subsection*{Your turn}
\begin{itemize}
\raggedright
  \item \emph{Say it:} \emph{vaśapaḍisiko}
  \item \emph{Say it:} \emph{vaśapaḍisiko}, once more
  \item \emph{Say it:} \emph{daṅgekōraru kōṭeyannu vaśapaḍisikoṇḍaru.} — The rebels captured the fort.
  \item \emph{Recall:} say the Kannada for a slave, then the Kannada for to set free, then say \emph{vaśapaḍisiko} again
\end{itemize}

\begin{tcolorbox}[breakable,colback=teal!4,colframe=teal!35!black,title={Before you move on}]
What does \kn{ವಶಪಡಿಸಿಕೊ} mean? (\textquotedblleft{}To capture; to take control of\textquotedblright{}.) Say it once more.
\end{tcolorbox}
